



%%% Here is the class with everything we need, if you don't want, you don't have to take a look at it, as it will probably cover your needs by far. These three options are used after the class article, which this template is based in.
% Twoside will act as an special parameter, because if the document is twosided indeed, headers and footers will be adapted in consequence.
% I recently added the language parameter to change it more easily, so the label of all theorems are redefined accordingly, depending on which is used. Disclaimer: only available for Catalan, English and Spanish. If no language is selected, English is enforced.


% \documentclass[a4paper,12pt,oneside,catalan]{all-in-one} %% ONESIDE
\documentclass[a4paper,12pt,twoside,portuguese]{all-in-one} %% TWOSIDE

% Here is the font I usually use, it can be changed of course.
% In that case, these next two lines have to be deleted.
% \usepackage{ebgaramond}
% Equations have the same font.
% \usepackage[cmintegrals,cmbraces,ebgaramond]{newtxmath}

\doctitle{SME 0340 - EDO}
\docsubtitle{Nome do Modelo Matemático}

\author{Andre V. V. Codorniz \\ \texttt{ID} \and Filipe S. Lopes \\ \texttt{ID} \and Heloísa Pazeti \\ \texttt{14577991} \and Lucas M. Farias \\ \texttt{ID} \and Miguel R. Fonseca
 \\ \texttt{14682196} \and Renan C. M. Soares \\ \texttt{ID}}

% A text that goes into the footer. You can even include a logo. It can be left blank, too.
%% Option 1: \footext{\includegraphics[width=2em]{titlepage.png}}
%% Option 2: \footext{\copyright\ \the\year by Cookie Monster}
% What I usually use: 
\footext{NOME DO MODELO MATEMÁTICO}

% Don't comment the command below if it's not of much inconvenience, please!
\greetings
% Licensing on every page? Comment if you don't want it.
% \licensing

% Last things
% This is my favourite config for bibliography. Each one of the references can be briefly explained with the field "note".
\RequirePackage[style=english]{csquotes}
\RequirePackage[
backend=biber,
style=alphabetic,
sorting=ynt
]{biblatex}
\newcommand{\familynameformat}[1]{\MakeUppercase{#1}}
\AtBeginBibliography{%
  \renewcommand{\mkbibnamefamily}{\familynameformat}%
}
\renewbibmacro*{begentry}{%
  \iffieldundef{note}
    {\undef\bbxnote}
    {\savefield{note}{\bbxnote}%
     \clearfield{note}}}
\renewbibmacro*{finentry}{%
  \restorefield{note}{\bbxnote}%
  \iffieldundef{note}
    {\finentry}
    {\setunit{\finentrypunct\par\vspace{\bibitemsep}\nobreak}
     \textit{\printfield{note}}%
     \finentry}}
\let\familynameformat=\textsc
\nocite{*}

\addbibresource{bibl.bib}


% Hyperref always required last one.
\RequirePackage{hyperref}
\makeatletter
\hypersetup{%
  hidelinks,
  pdfstartview=Fit,%
  pdfmenubar=true,%
  pdftoolbar=true,%
  bookmarksopen=false,%
  pdftitle={\@docsubtitle},%
  pdfauthor={\templateauthor},%
  pdfsubject={\@doctitle},%
  pdflang={\languagename},%
  pdfkeywords={mathematics},%
  pdfproducer={pdfTeX}}
\makeatother

% Setting up the titlepage.
\makeatletter%
\title{{\large\textit{\@doctitle}}\\[0.5cm]{\Huge\color{gray}\textsc{\@docsubtitle}}}%
\makeatother%
\date{\today}

% Dummy text
\usepackage{lipsum}

\begin{document}

\thispagestyle{empty}

\begin{figure}
  \centering
  \includegraphics[scale=0.2]{images/titlepage.png} % logo no topo da página
  \vspace*{-0.3cm}
\end{figure}

% \begin{figure}
%   \centering
%   \includegraphics[scale=0.08]{images/uni-kassel-logo.png}
%   \hspace*{2cm}
%   \includegraphics[scale=0.7]{images/mrt-logo.png}
%   \vspace*{-0.3cm}
% \end{figure}

\begin{center}
    \vspace*{-2cm}
\MakeUppercase{\Large{Universidade de São Paulo}} \\
\MakeUppercase{\Large{Escola de Engenharia de São Carlos}}\\
{\large \rm \textbf {Circuitos Elétricos - SEL0602} \linebreak} % trocar aqui de acordo com a matéria
\end{center}

\baselineskip 30pt

\vspace*{0.3cm}

\begin{center}
    {\LARGE \bfseries Atividade PSPICE - 1} % trocar aqui o título do trabalho
\end{center}

\vspace*{1cm}

\setcounter{footnote}{1}

\renewcommand{\thefootnote}{\fnsymbol{footnote}}
\begin{center}
{
   
    \sc  André Vargas Villalba Codorniz - 14558436 \\ % adicionar, deletar ou editar os nomes dos membros
   
}
\vspace*{1cm}

\center{\rm  PROFESSOR:} 
\vspace*{-.8cm}
\center{
    \sc Mário Oleskovicz \\ % trocar aqui o nome do professor
    \sc    % trocar aqui o nome do professor
} 
\end{center}

\baselineskip 17pt

\vspace*{1.5cm}
\begin{center}
    {{\bf São Carlos, Abril de 2024}} % trocar aqui o mês e ano do trabalho
\end{center}

\vspace*{.05cm}

\renewcommand{\thefootnote}{\arabic{footnote}}

\setcounter{footnote}{1}

\pagebreak

\baselineskip 19pt
\newpage

% Interlineate, it can be changed.
\setlength{\baselineskip}{.70cm}

% \begin{titlepage}
% \maketitle\vfill
% %\centering\includegraphics[width=20em]{titlepage.png}
% %\doclicenseThis
% % Very important, regarding page numbering. Keep in mind differences between oneside and twoside
% \thispagestyle{empty}
% \end{titlepage}

% Table of contents
\thispagestyle{plain}
\renewcommand*\contentsname{Sumário}
\tableofcontents
\newpage

\pagestyle{\defaultsettings}

\section{Introdução}

- Apresentação do Modelo Matemático.
- Descreva detalhadamente o modelo que será estudado e a sua importância.
\newpage

\section{Desenvolvimento}

- Desenvolvimento do modelo.

\subsection{Seção A}

- Resolução do modelo utilizando autovalores e autovetores.
- Descreva todos os passos da resolução do sistema.

\begin{equation}
    \begin{cases}
      x_1'(t)=-\frac{f_1}{V_1}x_1(t)-\frac{f_3}{V_3}x_3(t)+f(t)\\
      x_2'(t)=\frac{f_1}{V_1}x_1(t)-\frac{f_2}{V_2}x_2(t)\\
      x_3'(t)=\frac{f_2}{V_2}x_2(t)+\frac{f_3}{V_3}x_3(t)
    \end{cases}
\end{equation}

Em forma de matriz:
\begin{center}

    
\begin{equation}
\renewcommand{\arraystretch}{1.8}
\underbrace{\begin{pmatrix}
    x_1'(t)\\
    x_2'(t)\\
    x_3'(t)
\end{pmatrix}}_\text{$X'(t)$}
=
\underbrace{\begin{pmatrix}
-\frac{f_1}{V_1} & 0 & -\frac{f_3}{V_3}\\
\frac{f_1}{V_1} & -\frac{f_2}{V_2} & 0\\
0 & \frac{f_2}{V_2} & \frac{f_3}{V_3}
\end{pmatrix}}_\text{$A(t)$}
\cdot
\underbrace{\begin{pmatrix}
x_1(t)\\
x_2(t)\\
x_3(t)
\end{pmatrix}}_\text{$X(t)$}
+
\underbrace{\begin{pmatrix}
f(t)\\
0\\
0
\end{pmatrix}}_\text{$f(t)$}
\end{equation}
\end{center}



\subsection{Seção B}

- Análise gráfica do modelo.
- Faça uma análise qualitativa do modelo através de gráficos. 
- Apresente todos os cálculos para esboçar o gráfico do modelo. 
- Utilize a teoria de derivadas para estudar o gráfico do modelo


OBS: Se precisar usar alguma prova, corolário, ..., usar os templates abaixo.

% Theorem environments
\begin{theorem}
THEOREMS
\end{theorem}
\begin{proof}
\end{proof}
\begin{corollary}
COROLLARY
\end{corollary}
\begin{definition}
DEFINITION
\end{definition}
\begin{remark}
REMARK
\end{remark}
\begin{exercise}
EXERCISE
\end{exercise}
\begin{proof}[Solução]
SOLUTION
\end{proof}
\clearpage
\newpage

% Appendixes
\appendix
\section{Apêndice}

\newpage

% Bibliography
\clearpage
\pagestyle{\auxsettings}
\printbibliography[heading=bibintoc]

\end{document}

